\documentclass[11pt,a4paper]{article}
\usepackage[margin=1in]{geometry}
\usepackage[T1]{fontenc}
\usepackage{titlesec}
\usepackage{enumitem}
\usepackage{booktabs}
\usepackage{longtable}
\usepackage{xcolor}
\usepackage{fancyhdr}
\usepackage{hyperref}
\usepackage{parskip}
\usepackage{listings}

\hypersetup{
  colorlinks=true,
  linkcolor=blue!60!black,
  urlcolor=blue!60!black,
  citecolor=blue!60!black
}

\definecolor{codebg}{RGB}{244,244,244}
\lstset{
  backgroundcolor=\color{codebg},
  basicstyle=\ttfamily\small,
  breaklines=true,
  frame=single,
  columns=fullflexible,
  showstringspaces=false
}

\pagestyle{fancy}
\fancyhf{}
\fancyhead[L]{PetSocial --- Security and Privacy}
\fancyfoot[C]{\thepage}
\renewcommand{\headrulewidth}{0.4pt}

\title{PetSocial\\\large Security and Privacy}
\author{PetSocial Development Documentation}
\date{\today}

\begin{document}
\maketitle
\thispagestyle{fancy}
\tableofcontents

\section{Authentication}
Authentication is delegated entirely to Supabase Auth. Sessions are stored as
cookies and read through three context-specific Supabase client constructors (browser,
server, middleware --- see the System Architecture document). A Next.js middleware
runs on every request purely to refresh a session that is close to expiring, so that a
Server Component never has to handle an expired-but-refreshable token itself.

No page or Server Action ever trusts a user identifier passed from the client. Every
authenticated action re-derives the caller's identity server-side by calling
\texttt{requireUser()} (\texttt{src/lib/auth.ts}), which reads the session from cookies
and redirects to \texttt{/login} if none is present:
\begin{lstlisting}
export async function requireUser() {
  const supabase = await createClient();
  const { data: { user } } = await supabase.auth.getUser();
  if (!user) redirect("/login");
  return user;
}
\end{lstlisting}

\section{Authorisation}
Beyond ``is someone logged in,'' every action that mutates a specific resource
independently verifies that the resource belongs to the caller before proceeding ---
ownership is never assumed from a form field. For example, updating a pet's profile
looks up the pet scoped to both its id \emph{and} the caller's id in a single query:
\begin{lstlisting}
const pet = await prisma.pet.findFirst({
  where: { id: petId, ownerId: user.id },
});
if (!pet) return { error: "..." };
\end{lstlisting}
The same pattern is used for editing a pet, updating an adoption or marketplace
listing's status, cancelling a vet appointment, and resolving a lost-and-found report.

\section{Input Validation}
Every Server Action that accepts user input parses it through a Zod schema before
touching the database, with explicit length limits on every text field (documented
per-action in the API Reference). This serves two purposes: rejecting malformed input
with a clear error message, and preventing unbounded-length text from being written to
the database, which both bounds storage growth and reduces the attack surface for
downstream rendering.

\section{Rate Limiting}
Every mutating action is rate-limited per authenticated user (or, for the unauthenticated
sign-up action, per client IP address), using an in-process, in-memory sliding-window
counter (\texttt{src/lib/rate-limit.ts}). This is a deliberate simplification: it is
correct and effective for a single running server instance, but has two consequences
worth tracking as the deployment scales:
\begin{itemize}[leftmargin=*]
  \item Multiple concurrent serverless instances each hold independent counters, so the
        effective limit scales with instance count rather than staying fixed.
  \item Counters reset on every cold start or redeploy.
\end{itemize}
Before relying on these limits under real adversarial load, replace the backing store
with a shared one (e.g.\ Upstash Redis) behind the same function signature.

\section{The Anonymity Model}
An owner account and each individual pet profile carry independent
\texttt{isAnonymous} flags. Setting either hides the owner's real display name and
avatar from \emph{public display} --- it never removes or anonymises the underlying
account record, and it never prevents moderation or direct messaging, both of which
continue to resolve to the real account internally. The resolution rule is centralised
in one function, reused by every page that renders authorship, rather than
reimplemented per page:
\begin{lstlisting}
export function isHidden(entity, owner) {
  return Boolean(entity?.isAnonymous) || owner.isAnonymous;
}
\end{lstlisting}
Centralising this logic matters for correctness: before this refactor, several pages
independently implemented the same check with slightly different logic, which is
exactly the kind of drift that produces a page which forgets to hide an anonymous
owner's name.

\section{Owner-Only Data}
A pet's health log is the one class of data in the product that is private rather than
merely de-identified. It is never fetched at all when the requesting user is not the
pet's own owner --- the guard is on the database query itself, not on whether the
already-fetched data is rendered:
\begin{lstlisting}
const healthLogs = isOwner
  ? await prisma.healthLog.findMany({ where: { petId: pet.id } })
  : [];
\end{lstlisting}
This means a health log can never leak through a client-side rendering bug, a cache, or
a debugging tool inspecting the page's data, because it is simply never present in the
response in the first place for a non-owner viewer.

\section{File Upload Security}
User photo uploads go directly from the browser to Supabase Storage (not proxied
through the Next.js server), into a single public bucket, \texttt{post-media}, with two
Row-Level Security policies:
\begin{itemize}[leftmargin=*]
  \item Only authenticated users may insert new objects.
  \item Anyone may read any object (the bucket is intentionally public, since post
        photos and avatars are public content by nature).
\end{itemize}
There is currently no server-side validation of uploaded file type, size, or content
beyond what the browser's \texttt{accept="image/*"} hint encourages; a determined
authenticated user could upload a non-image file to the bucket. This is a known gap,
acceptable for the current product stage but worth closing (server-side content-type
and size checks) before handling untrusted public traffic at scale.

\section{Secrets Management}
\begin{itemize}[leftmargin=*]
  \item Database connection strings and Supabase API keys live only in \texttt{.env}
        and \texttt{.env.local}, both git-ignored; only \texttt{.env.example} (with
        placeholder values) is committed.
  \item Database credentials should be rotated any time they are known to have been
        shared outside of the Supabase dashboard or a properly access-controlled
        secrets manager --- for instance, after being pasted into a chat transcript,
        a support ticket, or any other non-durable channel, even if that channel is
        believed private.
  \item The publishable Supabase key (\texttt{NEXT\_PUBLIC\_SUPABASE\_PUBLISHABLE\_KEY})
        is, by design, safe to expose to the browser; it authorises only what the
        project's Row-Level Security policies allow.
\end{itemize}

\section{Known Limitations}
\begin{itemize}[leftmargin=*]
  \item \textbf{No automated test suite.} Every feature described in this document set
        was verified through disposable, manually-run scripts exercising the real
        application and database during development, not through a committed,
        repeatable test suite. Regressions currently rely on manual re-verification.
  \item \textbf{In-memory rate limiting}, as discussed in \S~4.
  \item \textbf{No server-side file validation}, as discussed in \S~7.
  \item \texttt{Owner.isPrivate} exists in the schema but is not enforced anywhere ---
        it does not currently make an account private in any functional sense.
\end{itemize}

\end{document}
